\documentclass[../../../include/open-logic-section]{subfiles}

\begin{document}

\olfileid{sth}{replacement}{ref}
\olsection{ప్రతిస్థాపన, ప్రతిబింబనం}

\stagesinex{} ద్వారా ప్రతిస్థాపనను సమర్థించే మన చివరి ప్రయత్నం ఒక గంభీరమైన, అందమైన ఫలితంతో మొదలవుతుంది:\footnote{గుర్తుంచుకోండి: వేరుగా స్పష్టంగా చెప్పనంతవరకు అన్ని సూత్రాల్లో పరామితులు ఉండవచ్చు.}
%In this, I will use overlining, such as $x_1, \ldots, x_n$, to
%abbreviate ``$x_1, \ldots, x_n$'':
\begin{thm}[ప్రతిబింబన పథకం]\ollabel{reflectionschema}
ఏ సూత్రం $\phi$కైనా:
\[
\forall \alpha \exists \beta > \alpha (\forall x_1 \ldots, x_n \in
V_\beta)(\phi(x_1, \ldots, x_n) \liff \phi^{V_\beta}(x_1, \ldots, x_n))
\]
\end{thm}
\noindent
\olref[sth][replacement][strength]{formularelativization}లోలాగే, $\phi^{V_\beta}$ అనేది $\phi$లోని ప్రతి పరిమాణకారకాన్ని~$V_\beta$ సమితికి పరిమితం చేస్తే వచ్చే సూత్రం. కనుక సహజార్థంలో ప్రతిబింబనం ఇలా చెబుతుంది: మొత్తం స్థాయిక్రమంలో $\phi$ నిజమైతే, స్థాయిక్రమపు అపరిమితంగా అనేక \emph{ఆరంభ ఖండాల్లో} కూడా $\phi$ నిజమే.

\citet{Montague1961}, \citet{Levy1960} చూపినదేమిటంటే, $\Z$ పరంగా ప్రతిస్థాపన, ప్రతిబింబనం (తగిన రూపాల్లో) సమానార్థకాలు; కాబట్టి వాటిలో ఏదైనా చేర్చితే~$\ZF$ వస్తుంది. (ఈ ఫలితాలను \olref[sth][replacement][refproofs]{sec}లో నిరూపిస్తాం.) ఈ సమానార్థకత్వం వల్ల \stagesinex{} ఆధారంగా ప్రతిబింబనను, ప్రతిస్థాపనను ఇలా సమర్థించవచ్చని ఆశించవచ్చు: \stagesinex{} ఉన్నప్పుడు స్థాయిక్రమం చాలా, చాలా ఎత్తుగా ఉండాలి; దాని గురించి మనం చెప్పగల ఏదీ దాని ఎత్తుకు హద్దు పెట్టలేనంత ఎత్తుగా. దీనిని, ఏ వాక్యం~$\phi$ మొత్తం స్థాయిక్రమంలో నిజమైతే అది అపరిమితంగా అనేక ఆరంభ ఖండాల్లోనూ నిజమనే ఆలోచనగా అర్థం చేసుకోవచ్చు. అదే ప్రతిబింబనం.

ఇదీ ప్రతిస్థాపనకు అంతర్గత సమర్థన ఇచ్చే నిజంగా ఆశాజనకమైన ప్రయత్నంలా కనిపిస్తుంది. కానీ ఇక్కడ దీని గురించి చెప్పాల్సినది చాలా ఎక్కువ. ఇది విజయవంతమవుతుందో లేదో ఇప్పుడు మీరే నిర్ణయించాలి.\footnote{అయితే తరువాత \citet[95--100]{Incurvati2020} చదవవచ్చు.}

\end{document}
